\documentclass[11pt]{article}
% Shared print layout. The web's provisional design and MathJax are independent.
\usepackage[utf8]{inputenc}
\usepackage[T1]{fontenc}
\usepackage[spanish,es-nodecimaldot]{babel}
\usepackage{lmodern}
\usepackage{amsmath,amssymb}
\usepackage{graphicx}
\usepackage{geometry}
\geometry{a4paper,left=20mm,right=20mm,top=25mm,bottom=22mm,headheight=16pt,headsep=8mm}
\usepackage{fancyhdr,lastpage}
\usepackage{enumitem}
\usepackage{xcolor}
\usepackage{pgfplots}
\pgfplotsset{compat=1.18}
\usepgfplotslibrary{fillbetween}
\usepackage[hidelinks,pdfencoding=auto]{hyperref}
\hypersetup{pdflang=es-ES,pdfauthor={},pdfsubject={Matemáticas Aplicadas a las Ciencias Sociales II}}
\pagestyle{fancy}
\fancyhf{}
\fancyhead[L]{\small Matemáticas Aplicadas a las Ciencias Sociales II}
\fancyhead[R]{\small 2.\textsuperscript{o} Bachillerato}
\fancyfoot[L]{\footnotesize Límites y derivadas | Edición revisada 1.0}
\fancyfoot[R]{\footnotesize Página \thepage\ de \pageref{LastPage}}
\renewcommand{\headrulewidth}{0.4pt}
\setlength{\parindent}{0pt}
\setlength{\parskip}{0.5em}
\setlist{itemsep=0.3em,topsep=0.3em}
\newcommand{\apartado}[1]{\subsection*{#1}}
\pgfplotsset{
  revision/.style={
    width=0.94\linewidth,height=7.5cm,
    axis lines=middle,grid=major,grid style={gray!18},
    tick label style={font=\small},label style={font=\small},
    xlabel={$x$},ylabel={$y$},samples=160,
    legend style={at={(0.5,-0.16)},anchor=north,draw=none,font=\small},
    clip=true
  }
}

\hypersetup{pdftitle={Límites y derivadas - Soluciones revisadas}}
\begin{document}
\section*{Ejercicio 1: continuidad y derivabilidad}
Solucionario elaborado en esta revisión a partir de F008. El original solo contiene comentarios parciales del ejercicio 4.

En la unión $0$, $f(0)=2$. La rama racional tiene límite derecho $-b$, de modo que continuidad exige $b=-2$.
\[f'(x)=\begin{cases}2x+a,&x<0,\\\dfrac{-1-b}{(x-1)^2},&x>0,\ x\neq1.\end{cases}\]
Con $b=-2$, las derivadas laterales en $0$ son $a$ y $1$. Igualándolas:
\[\boxed{a=1,\qquad b=-2}.\]
Para estos parámetros, los cocientes incrementales respecto a $0$ son $h+1$ a la izquierda y $1/(1-h)$ a la derecha; ambos tienden a 1. El punto $x=1$ queda fuera del dominio, pero no afecta al estudio en $0$.
\section*{Ejercicio 1(b,c): tangente y normal}
Se utilizan los parámetros obtenidos en (a). En $2$ rige la rama racional:
\[f(2)=0,\qquad f'(2)=1.\]
La recta tangente pasa por $(2,0)$ y tiene pendiente 1:
\[\boxed{y=x-2}.\]
La normal pasa por el mismo punto y es perpendicular a ella. Su pendiente es $-1$:
\[\boxed{y=-x+2}.\]
\textbf{Comprobación:} ambas contienen $(2,0)$; el producto de sus pendientes es $-1$. El valor de $f'(2)$ confirma la pendiente tangente.
\clearpage
\section*{Ejercicio 2: derivadas (a--c)}
Se muestran las seis respuestas para consulta. En el examen se eligen cuatro y solo se corrigen las cuatro primeras según orden de lectura, a 0,6 puntos cada una.
\subsection*{(a) Cociente elevado al cubo}
Sea $u=(3x^2-5)/(2x+3)$. Entonces
\[u'=\frac{6x(2x+3)-2(3x^2-5)}{(2x+3)^2}=\frac{6x^2+18x+10}{(2x+3)^2}.\]
Aplicando la regla de la cadena,
\[\boxed{f'=\frac{3(3x^2-5)^2(6x^2+18x+10)}{(2x+3)^4}},\qquad x\neq-3/2.\]
\subsection*{(b) Cuadrado del seno}
\[\boxed{f'=2\sin x\cos x=\sin(2x)},\qquad x\in\mathbb R.\]
Se deriva primero el cuadrado y después su función interior.
\subsection*{(c) Cociente con denominador compuesto}
\[f'=\frac{3x^2(x^2+4)^2-(x^3-4)\,4x(x^2+4)}{(x^2+4)^4}.\]
Cancelando un factor positivo $x^2+4$,
\[\boxed{f'=\frac{-x^4+12x^2+16x}{(x^2+4)^3}},\qquad x\in\mathbb R.\]
El denominador no se anula para ningún valor real.
\clearpage
\section*{Ejercicio 2: derivadas (d--f)}
\subsection*{(d) Exponencial compuesta}
\[\boxed{f'=(6x-1)e^{3x^2-x+2}},\qquad x\in\mathbb R.\]
Se conserva la exponencial y se multiplica por la derivada del exponente.
\subsection*{(e) Coseno compuesto}
\[\boxed{f'=-(2x+2)\sin(x^2+2x)},\qquad x\in\mathbb R.\]
El signo negativo procede de la derivada del coseno.
\subsection*{(f) Cociente con raíz y logaritmo}
El denominador debe ser real y no nulo: $1+2\ln x>0$. Por tanto, el dominio es $x>e^{-1/2}$.
\[f=x^2(1+2\ln x)^{-1/2}.\]
Por producto y cadena,
\begin{align*}
f'&=2x(1+2\ln x)^{-1/2}-x(1+2\ln x)^{-3/2}\\
&=\boxed{\frac{x(1+4\ln x)}{(1+2\ln x)^{3/2}}}.
\end{align*}
No basta con exigir $x>0$: la raíz del denominador impone la condición más restrictiva indicada.
\clearpage
\section*{Ejercicio 3: límites}
Se muestran los seis. Se corrigen solo los cuatro primeros elegidos, según orden de lectura.
\subsection*{(a) Racionalización}
\[\frac{1-\sqrt{1-x^2}}{x^2}=\frac1{1+\sqrt{1-x^2}}\ \longrightarrow\ \boxed{\frac12}.\]
\subsection*{(b) Exponencial frente a polinomio}
La exponencial crece más rápido que cualquier potencia de $x$:
\[3^x-x^5=3^x\left(1-\frac{x^5}{3^x}\right)\longrightarrow\boxed{+\infty}.\]
\subsection*{(c) División por $x>0$}
\[\frac{5x+7}{\sqrt{2x^2+3}}=\frac{5+7/x}{\sqrt{2+3/x^2}}\longrightarrow\boxed{\frac5{\sqrt2}}.\]
\subsection*{(d) Potencia indeterminada}
Para $x$ grande la base es positiva. Si $L_x=(1-1/x^2)^x$,
\[\ln L_x=x\ln(1-1/x^2)\sim-1/x\longrightarrow0,\qquad\boxed{L_x\longrightarrow1}.\]
\subsection*{(e) Denominador común}
\[\frac7{x-4}-\frac{4x+12}{x(x-4)}=\frac{3x-12}{x(x-4)}=\frac3x\longrightarrow\boxed{\frac34}.\]
La igualdad simplificada se usa para $x\neq4$, suficiente para calcular el límite.
\subsection*{(f) Límites laterales distintos}
$x^2-3x+2=(x-1)(x-2)$. Por la izquierda de 1, ambos factores son negativos; por la derecha tienen signos opuestos:
\[\lim_{x\to1^-}f=+\infty,\qquad\lim_{x\to1^+}f=-\infty.\]
El límite bilateral \textbf{no existe}, ni siquiera como un único infinito con signo.
\clearpage
\section*{Ejercicio 4: distancia del cohete}
El tiempo físico cumple $t\geq0$. Al comenzar la fase final,
\[D(0)=100\text{ miles de kilómetros}=\boxed{100\,000\text{ km}}.\]
Después de dos horas,
\[D(2)=\frac{100}{5}=20\text{ miles de kilómetros}=\boxed{20\,000\text{ km}}.\]
Como $t^2+1\to+\infty$,
\[\boxed{\lim_{t\to+\infty}D(t)=0}.\]
\textbf{Corrección M14.} Para todo tiempo finito $t\geq0$,
\[D(t)=\frac{100}{t^2+1}>0.\]
El modelo se aproxima a la zona de aterrizaje, pero \textbf{nunca alcanza distancia cero en un tiempo finito}. No permite afirmar que el cohete aterrice ni calcular un instante de llegada. El comentario del original que deduce aterrizaje a partir del límite no es correcto.
\begin{center}\begin{tikzpicture}\begin{axis}[revision,height=7cm,xmin=0,xmax=12,ymin=0,ymax=110,xtick={0,2,4,8,12},ytick={0,20,60,100},xlabel={Horas ($t$)},ylabel={Miles de kilómetros}]
\addplot[thick,domain=0:12] {100/(x^2+1)};
\addplot[only marks,mark=*] coordinates {(0,100) (2,20)};
\end{axis}\end{tikzpicture}\end{center}
La gráfica desciende y se aproxima al eje; no se añade un punto de impacto ficticio.
\clearpage
\section*{Anexo no evaluable: derivadas comentadas}
Se conservan las seis alternativas comentadas del original, sin añadir puntos ni modificar las opciones activas. Se normalizan los delimitadores LaTeX dañados de la primera expresión.
\subsection*{(a) $f=((x^2+2)/(4x+2))^2$}
\[\boxed{f'=\frac{2(x^2+2)(4x^2+4x-8)}{(4x+2)^3}},\qquad x\neq-1/2.\]
\subsection*{(b) $f=\sin^2 x$}
\[\boxed{f'=2\sin x\cos x},\qquad x\in\mathbb R.\]
\subsection*{(c) $f=(2-x^2)/(x^2+1)^3$}
\[\boxed{f'=\frac{4x^3-14x}{(x^2+1)^4}},\qquad x\in\mathbb R.\]
\subsection*{(d) $f=e^{x^2+2x-1}$}
\[\boxed{f'=(2x+2)e^{x^2+2x-1}},\qquad x\in\mathbb R.\]
\subsection*{(e) $f=\cos(x^2+3x)$}
\[\boxed{f'=-(2x+3)\sin(x^2+3x)},\qquad x\in\mathbb R.\]
\subsection*{(f) $f=x/\sqrt{1+x}$}
\[\boxed{f'=\frac{x+2}{2(1+x)^{3/2}}},\qquad x>-1.\]
Los dominios se respetan antes de simplificar; no se usan valores donde el denominador o la raíz impiden definir la función.
\clearpage
\section*{Anexo no evaluable: límites comentados}
Se conservan los seis enunciados alternativos del original, con resultados comprobados.
\subsection*{(a) $\displaystyle\lim_{x\to+\infty}(e^x-1)/x^5$}
Domina la exponencial: el resultado es $\boxed{+\infty}$.
\subsection*{(b) $\displaystyle\lim_{x\to+\infty}(2^x-x^4)$}
Domina la exponencial: el resultado es $\boxed{+\infty}$.
\subsection*{(c) $\displaystyle\lim_{x\to+\infty}(3x-5)/\sqrt{x^2+1}$}
Al dividir por $x>0$ resulta $(3-5/x)/\sqrt{1+1/x^2}\to\boxed{3}$.
\subsection*{(d) $\displaystyle\lim_{x\to-\infty}(1+2/x)^{(x^2+2)/x}$}
La base es positiva para $x<-2$. El logaritmo de la expresión es
\[(x+2/x)\ln(1+2/x)\longrightarrow2,\qquad\text{resultado: }\boxed{e^2}.\]
\subsection*{(e) $\displaystyle\lim_{x\to3}\left(6/(x-3)-(3x+9)/(x^2-3x)\right)$}
\[\frac{6x-3x-9}{x(x-3)}=\frac3x\longrightarrow\boxed{1}.\]
\subsection*{(f) $\displaystyle\lim_{x\to2}1/(x^2-5x+6)$}
El denominador es $(x-2)(x-3)$. Los límites izquierdo y derecho son $+\infty$ y $-\infty$, respectivamente; el bilateral \textbf{no existe}.
\end{document}
